\begin{quote}\small Historical Round27 checkpoint. References below to planned, unexecuted or next work describe the state at that checkpoint; the integrated Round29 chapter supersedes that plan.\end{quote}
\section{Contributions and remaining proof obligations}
\begingroup\small
\begin{longtable}{@{}>{\raggedright\arraybackslash}p{24mm}>{\raggedright\arraybackslash}p{45mm}>{\raggedright\arraybackslash}p{34mm}>{\raggedright\arraybackslash}p{44mm}@{}}
\toprule Contribution & Earlier understanding $\to$ present change & Possible application & Limitation and next missing step \\
\midrule\endhead
AI1 fibers\workstar & Nontrivial endpoint formulas $\to$ exact classification of the combinations identified by each declared observation convention. & Remove redundant parameter fits; choose which calibration must be independent. & Surrogate identities are not finite-$q$ model equivalences. Prove and use actual finite-profile error disks. \\
\addlinespace
AI1 carrier\workstar & Demodulated slow signal $\to$ rank-two identification of carrier and slow rates with a residual $\eta$--$q$ fiber and explicit sampling aliases. & Plan phase-aware, nonaliased observation protocols. & Requires calibrated physical time/action/phase. No implemented clock or derivative transfer theorem. \\
\addlinespace
AI2 probes and support\workstar & One coherent face type $\to$ two actual multiplier sources with weights $q^4,q^5$, ten-factor common seed and recomputed complete collars. & Design observables sensitive to profile differences; prevent omitted-channel understatement. & Source-specific canonical family. The component is not the whole regional energy shell. \\
\addlinespace
AI2 full-scalar test\workstar & Retained-center differences $\to$ disjoint complete scalar disks for one frozen grid cell. Other cells and all simple ratio tests remain insufficient. & Reject one of two named model hypotheses when the full error allowance is met. & $T=2\times10^{48}\hbar/\alpha$ and roughly $10^{-26}$ scalar margins; no feasibility claim or continuous inverse. \\
\addlinespace
AG1 one update\workstar & Individual-source estimates $\to$ an entire nonlinear remainder bound and actual selected-source contraction below $2/3$ on the narrower interval. & Identify the terms, input denominator and norm costs an iteration must control. & Restricted $G+A$ source sector and one spent weight margin. Other original terms require explicit transport; no all-stage or complete-Hamiltonian contraction follows. \\
\bottomrule
\end{longtable}
\endgroup
The two-star notation marks workbench specialization, derivation or implementation, not established scientific priority. These results may be useful even when a desired certificate fails: a failed scalar separation or contraction upper bound identifies a limitation of the sufficient estimate. It does not prove equality, divergence or impossibility of the underlying dynamics.

\subsection{Why a percentage of solving Yang--Mills is not estimable}
The three requested investigations have final scoped gates: the research task count is $3/3$. That known denominator does not supply a percentage of a mathematical proof. The expert lenses and skeptic agree that no defensible numerical fraction of the Yang--Mills solution is identifiable from pages, loops, equations, network nodes or finite-model successes. The remaining obligations depend on each other and have no justified equal weights or known total amount of work.

The official problem asks for a nontrivial four-dimensional quantum Yang--Mills theory with the required axiomatic strength and a positive mass gap for every compact simple gauge group. The official status remains unsolved at the recorded source check. The workbench's outstanding obligations are:
\begin{center}\small
\begin{tabularx}{\textwidth}{@{}p{48mm}X@{}}
\toprule Obligation & Missing implication after Round27 \\
\midrule
Homogeneous control and matching & Convergent changed-source/changed-diagonal induction and a proved map on states, observables, domains and generators. \\
Uniform limiting estimates & Bounds surviving growing volume and regulator removal along a physically matched trajectory in the required topology. \\
Nontrivial continuum construction & The requested four-dimensional limiting theory and nontrivial local fields. \\
Axiomatic properties & Limiting locality, covariance, positivity/reconstruction and the required ultraviolet behavior. \\
Positive physical gap & A positive finite threshold in the actual constructed limiting theory, in fixed physical units. \\
Gauge-group coverage & The required compact simple groups beyond the workbench's $\mathrm{SU}(2)$ emphasis. \\
\bottomrule
\end{tabularx}
\end{center}
For a dimensionless cutoff operator $L(a)$ and energy convention $E(a)$, the identity
$\Delta_{\mathrm{phys}}(a)=E(a)\Delta_L(a)$ only converts units. A positive finite-cutoff $\Delta_L(a)$ proves no positive limiting physical gap without a justified scale, limiting construction and uniform estimate. An unknown completion percentage is not a claim of zero research value, nor an estimate of remaining effort.

The panel's next unexecuted priority is AG2: recover the original transformed Hamiltonian at the source checkpoint, inventory every remaining term in $E$, prove its stronger support-norm budget, and test whether its transported mixing cost fits inside the selected-source contraction margin. Only then should a changed-diagonal step be frozen with its own domain and remaining weight hypotheses. A cancellation-aware AI3 finite-value inverse remains an alternative observable branch. Neither is executed in this three-loop continuation.
